% Sorgente LaTeX canonico, importato e revisionato dal precedente sorgente Markdown.
\chapter{Calcolo differenziale}
\label{chap:10}

\begin{obiettivocapitolo}{calcolare derivate, tangenti, approssimazioni lineari e punti di non derivabilità}
\prioritaalta. Il dominio della funzione originaria prevale sempre sul dominio apparente della formula derivata.
\end{obiettivocapitolo}
\begin{proceduraesame}{derivata di un'espressione composta}
\begin{enumerate}
\item Identifica la funzione esterna e procedi dall'esterno verso l'interno.
\item Applica linearità, prodotto, quoziente e catena, scrivendo ogni fattore derivativo.
\item Semplifica solo dopo aver completato la derivazione.
\item Riporta le condizioni di esistenza della funzione e della derivata.
\end{enumerate}
\end{proceduraesame}
\begin{sceltametodo}{punti problematici}
Per funzioni a tratti o con valore assoluto, calcola le derivate laterali. Se sono finite e diverse c'è un punto angoloso; se sono infinite con segni opposti, una cuspide; se coincidono infinite, una tangente verticale. La derivabilità richiede continuità, ma la continuità non basta.
\end{sceltametodo}
\begin{proceduraesame}{retta tangente e linearizzazione}
Calcola $f(x_0)$ e $f'(x_0)$, quindi usa $y=f(x_0)+f'(x_0)(x-x_0)$. Per piccoli incrementi, $\Delta f\approx f'(x_0)\Delta x$; specifica sempre il punto di linearizzazione.
\end{proceduraesame}
\begin{controllofinale}{derivata}
Verifica con casi semplici, segno atteso, unità fisiche e derivazione inversa della primitiva quando applicabile.
\end{controllofinale}

La derivata misura il tasso di variazione locale di una funzione. Ogni regola va applicata soltanto nei punti nei quali la funzione originaria, le funzioni intermedie e la formula risultante sono definite. La notazione \(u'\) abbrevia \(du/dx\) quando \(u=u(x)\).

\section{Derivata, derivate laterali e retta tangente}\label{derivata-derivate-laterali-e-retta-tangente}

Siano \(f:D\to\mathbb R\) e \(x_0\in D\) punto di accumulazione di \(D\).

\[
\displaystyle
f'(x_0)=
\lim_{\substack{h\to0\\h\ne0,\ x_0+h\in D}}
\frac{f(x_0+h)-f(x_0)}{h}.
\]

Forma equivalente, ponendo \(x=x_0+h\):

\[
f'(x_0)
=\lim_{\substack{x\to x_0\\x\ne x_0,\ x\in D}}
\frac{f(x)-f(x_0)}{x-x_0}.
\]

Il numeratore è l'incremento \(\Delta f\) e il denominatore l'incremento \(\Delta x\); il loro rapporto è il rapporto incrementale.

\[
\displaystyle
f'_+(x_0)=
\lim_{\substack{h\to0^+\\x_0+h\in D}}
\frac{f(x_0+h)-f(x_0)}{h},
\qquad
f'_-(x_0)=
\lim_{\substack{h\to0^-\\x_0+h\in D}}
\frac{f(x_0+h)-f(x_0)}{h}.
\]

Se \(x_0\) è punto di accumulazione del dominio da entrambi i lati:

\[
\displaystyle
f'(x_0)\in\mathbb R
\quad\Longleftrightarrow\quad
f'_-(x_0)=f'_+(x_0)\in\mathbb R.
\]

Approssimazione lineare e retta tangente:

\[
\displaystyle
f(x_0+h)=f(x_0)+f'(x_0)h+o(h),
\]

\[
\displaystyle
y=f(x_0)+f'(x_0)(x-x_0).
\]

Retta normale al grafico nel punto \((x_0,f(x_0))\):

\[
y-f(x_0)=-\frac{1}{f'(x_0)}(x-x_0),
\qquad f'(x_0)\ne0,
\]

\[
x=x_0,
\qquad f'(x_0)=0.
\]

\[
\displaystyle
f\text{ derivabile in }x_0
\quad\Longrightarrow\quad
f\text{ continua in }x_0.
\]

Il viceversa non vale; caso standard:

\[
\displaystyle
f(x)=|x|,
\qquad
f'_-(0)=-1,
\qquad
f'_+(0)=1.
\]

Approfondimento: \href{https://ingegnerismo.it/matematica/derivata/}{derivata}.

\section{Punti di non derivabilità}\label{punti-di-non-derivabilituxe0}

\begin{longtable}[]{@{}
  >{\raggedright\arraybackslash}p{(\columnwidth - 4\tabcolsep) * \real{0.3333}}
  >{\raggedright\arraybackslash}p{(\columnwidth - 4\tabcolsep) * \real{0.3333}}
  >{\raggedright\arraybackslash}p{(\columnwidth - 4\tabcolsep) * \real{0.3333}}@{}}
\toprule\noalign{}
\begin{minipage}[b]{\linewidth}\raggedright
Caso
\end{minipage} & \begin{minipage}[b]{\linewidth}\raggedright
Derivate laterali
\end{minipage} & \begin{minipage}[b]{\linewidth}\raggedright
Modello in \(x_0=0\)
\end{minipage} \\
\midrule\noalign{}
\endhead
\bottomrule\noalign{}
\endlastfoot
Punto angoloso & \(f'_-,f'_+\in\mathbb R\), \(f'_-\ne f'_+\) & \(f(x)=\lvert x\rvert\) \\
Cuspide & \(f'_-,f'_+\) infiniti di segno opposto & \(f(x)=\sqrt{\lvert x\rvert}\) \\
Tangente verticale & \(f'_-=f'_+=\pm\infty\) & \(f(x)=\sqrt[3]{x}\) \\
Caso misto & un laterale è finito e l'altro infinito o inesistente & \(f(x)=\begin{cases}x,&x<0\\ \sqrt{x},&x\ge0\end{cases}\) \\
\end{longtable}

Approfondimento: \href{https://ingegnerismo.it/matematica/punti-di-non-derivabilita/}{punti di non derivabilità}.

\section{Differenziale e linearizzazione}\label{differenziale-e-linearizzazione}

Sia \(f:D\to\mathbb R\) derivabile in \(x_0\in D\cap D'\). Per \(h\to0\) con \(x_0+h\in D\):

\[
\displaystyle
df_{x_0}:h\longmapsto f'(x_0)h.
\]

\[
\displaystyle
f(x_0+h)-f(x_0)=f'(x_0)h+r(h),
\qquad
\lim_{h\to0}\frac{r(h)}{h}=0.
\]

\[
L_{x_0}(x)
=f(x_0)+df_{x_0}(x-x_0)
=f(x_0)+f'(x_0)(x-x_0).
\]

Con \(h=\Delta x\):

\[
\begin{aligned}
\Delta f&=f(x_0+h)-f(x_0),\\
df_{x_0}(h)&=f'(x_0)h,\\
r(h)&=\Delta f-df_{x_0}(h).
\end{aligned}
\]

Stime di primo ordine, valide per incrementi sufficientemente piccoli e con errore trascurabile rispetto a \(|\Delta x|\):

\[
\displaystyle
|\Delta f|\approx |f'(x_0)|\,|\Delta x|,
\]

\[
\displaystyle
\frac{|\Delta f|}{|f(x_0)|}
\approx
\left|\frac{f'(x_0)}{f(x_0)}\right|\,|\Delta x|,
\qquad f(x_0)\ne0.
\]

Approfondimento: \href{https://ingegnerismo.it/matematica/differenziale/}{differenziale e linearizzazione}.

Esercizi svolti: \href{https://ingegnerismo.it/matematica/differenziale-approssimazione-lineare-esercizi/}{differenziale e approssimazione lineare}.

\section{Derivate elementari}\label{derivate-elementari}

\begin{longtable}[]{@{}
  >{\raggedright\arraybackslash}p{(\columnwidth - 4\tabcolsep) * \real{0.3333}}
  >{\raggedright\arraybackslash}p{(\columnwidth - 4\tabcolsep) * \real{0.3333}}
  >{\raggedright\arraybackslash}p{(\columnwidth - 4\tabcolsep) * \real{0.3333}}@{}}
\toprule\noalign{}
\begin{minipage}[b]{\linewidth}\raggedright
Funzione
\end{minipage} & \begin{minipage}[b]{\linewidth}\raggedright
Derivata
\end{minipage} & \begin{minipage}[b]{\linewidth}\raggedright
Condizioni
\end{minipage} \\
\midrule\noalign{}
\endhead
\bottomrule\noalign{}
\endlastfoot
\(c\) & \(0\) & \(c\in\mathbb R\) \\
\(x\) & \(1\) & \(x\in\mathbb R\) \\
\(x^\alpha\) & \(\alpha x^{\alpha-1}\) & \(x>0\), \(\alpha\in\mathbb R\) \\
\(x^n\) & \(n x^{n-1}\) & \(x\in\mathbb R\), \(n\in\mathbb N_+\) \\
\(x^{-n}\) & \(-n x^{-n-1}\) & \(x\ne0\), \(n\in\mathbb N_+\) \\
\(\sqrt{x}\) & \(1/(2\sqrt{x})\) & \(x>0\) \\
\(\lvert x\rvert\) & \(\operatorname{sgn}x\) & \(x\ne0\); non derivabile in \(0\) \\
\(e^x\) & \(e^x\) & \(x\in\mathbb R\) \\
\(a^x\) & \(a^x\ln a\) & \(a>0\) \\
\(\ln x\) & \(1/x\) & \(x>0\) \\
\(\log_a x\) & \(1/(x\ln a)\) & \(x>0\), \(a>0\), \(a\ne1\) \\
\(\sin x\) & \(\cos x\) & \(x\in\mathbb R\) \\
\(\cos x\) & \(-\sin x\) & \(x\in\mathbb R\) \\
\(\tan x\) & \(1/\cos^2x=1+\tan^2x\) & \(x\ne\pi/2+k\pi\) \\
\(\cot x\) & \(-1/\sin^2x\) & \(x\ne k\pi\) \\
\(\sec x\) & \(\sec x\tan x\) & \(\cos x\ne0\) \\
\(\csc x\) & \(-\csc x\cot x\) & \(\sin x\ne0\) \\
\end{longtable}

\section{Derivate delle funzioni inverse}\label{derivate-delle-funzioni-inverse}

\begin{longtable}[]{@{}
  >{\raggedright\arraybackslash}p{(\columnwidth - 4\tabcolsep) * \real{0.3333}}
  >{\raggedright\arraybackslash}p{(\columnwidth - 4\tabcolsep) * \real{0.3333}}
  >{\raggedright\arraybackslash}p{(\columnwidth - 4\tabcolsep) * \real{0.3333}}@{}}
\toprule\noalign{}
\begin{minipage}[b]{\linewidth}\raggedright
Funzione
\end{minipage} & \begin{minipage}[b]{\linewidth}\raggedright
Derivata
\end{minipage} & \begin{minipage}[b]{\linewidth}\raggedright
Dominio della derivata
\end{minipage} \\
\midrule\noalign{}
\endhead
\bottomrule\noalign{}
\endlastfoot
\(\arcsin x\) & \(1/\sqrt{1-x^2}\) & \(\lvert x\rvert<1\) \\
\(\arccos x\) & \(-1/\sqrt{1-x^2}\) & \(\lvert x\rvert<1\) \\
\(\arctan x\) & \(1/(1+x^2)\) & \(x\in\mathbb R\) \\
\(\operatorname{arccot}x\) & \(-1/(1+x^2)\) & \(x\in\mathbb R\), con valori in \((0,\pi)\) \\
\end{longtable}

\section{Derivate iperboliche}\label{derivate-iperboliche}

\begin{longtable}[]{@{}lll@{}}
\toprule\noalign{}
Funzione & Derivata & Dominio della derivata \\
\midrule\noalign{}
\endhead
\bottomrule\noalign{}
\endlastfoot
\(\sinh x\) & \(\cosh x\) & \(x\in\mathbb R\) \\
\(\cosh x\) & \(\sinh x\) & \(x\in\mathbb R\) \\
\(\tanh x\) & \(1/\cosh^2x\) & \(x\in\mathbb R\) \\
\(\coth x\) & \(-1/\sinh^2x\) & \(x\ne0\) \\
\(\operatorname{arsinh}x\) & \(1/\sqrt{1+x^2}\) & \(x\in\mathbb R\) \\
\(\operatorname{arcosh}x\) & \(1/\sqrt{x^2-1}\) & \(x>1\) \\
\(\operatorname{artanh}x\) & \(1/(1-x^2)\) & \(\lvert x\rvert<1\) \\
\(\operatorname{arcoth}x\) & \(1/(1-x^2)\) & \(\lvert x\rvert>1\) \\
\end{longtable}

Approfondimenti: \href{https://ingegnerismo.it/matematica/derivata/}{derivata} e \href{https://ingegnerismo.it/matematica/derivata-ordine-superiore/}{derivate di ordine superiore}.

\section{Regole di derivazione}\label{regole-di-derivazione}

\[
\begin{aligned}
(\alpha f+\beta g)'&=\alpha f'+\beta g',\\
(fg)'&=f'g+fg',\\
\left(\frac1f\right)'&=-\frac{f'}{f^2},\qquad f\ne0,\\
\left(\frac fg\right)'&=\frac{f'g-fg'}{g^2},\qquad g\ne0,\\
(f\circ g)'(x)&=f'(g(x))g'(x),\\
(|f|)'&=\operatorname{sgn}(f)f',\qquad f\ne0.
\end{aligned}
\]

Nei punti in cui \(f(x)=0\), la derivabilità di \(|f|\) va controllata direttamente: la formula con \(\operatorname{sgn}(f)\) non decide il caso.

Se \(f\) è continua e strettamente monotona su un intervallo \(I\), è derivabile in \(x_0\in\operatorname{int}I\) e \(f'(x_0)\ne0\), posto \(y_0=f(x_0)\):

\[
\displaystyle
(f^{-1})'(y_0)=\frac{1}{f'(x_0)}.
\]

In forma operativa, per ogni \(y\) nel tratto nel quale valgono le stesse ipotesi:

\[
(f^{-1})'(y)
=\frac{1}{f'\bigl(f^{-1}(y)\bigr)}.
\]

Derivate successive:

\[
\displaystyle
f^{(0)}=f,
\qquad
f^{(n)}=(f^{(n-1)})'.
\]

Formula di Leibniz:

\[
\displaystyle
(fg)^{(n)}
=
\sum_{k=0}^{n}\binom nk f^{(k)}g^{(n-k)}.
\]

Derivata logaritmica:

\[
\displaystyle
(\ln|f(x)|)'=\frac{f'(x)}{f(x)},
\qquad f(x)\ne0.
\]

Potenza a base ed esponente variabili:

\[
\displaystyle
\bigl(u(x)^{v(x)}\bigr)'
=
u(x)^{v(x)}
\left[
v'(x)\ln u(x)+v(x)\frac{u'(x)}{u(x)}
\right],
\qquad u(x)>0.
\]

Approfondimenti: \href{https://ingegnerismo.it/matematica/regole-di-derivazione/}{regole di derivazione}, \href{https://ingegnerismo.it/matematica/regola-della-catena/}{regola della catena} e \href{https://ingegnerismo.it/matematica/derivata-logaritmica/}{derivata logaritmica}.

\section{Tabella delle derivate composte}\label{tabella-delle-derivate-composte}

Per \(u=u(x)\) derivabile, nei rispettivi domini:

\begin{longtable}[]{@{}
  >{\raggedright\arraybackslash}p{(\columnwidth - 4\tabcolsep) * \real{0.3333}}
  >{\raggedright\arraybackslash}p{(\columnwidth - 4\tabcolsep) * \real{0.3333}}
  >{\raggedright\arraybackslash}p{(\columnwidth - 4\tabcolsep) * \real{0.3333}}@{}}
\toprule\noalign{}
\begin{minipage}[b]{\linewidth}\raggedright
Funzione composta
\end{minipage} & \begin{minipage}[b]{\linewidth}\raggedright
Derivata
\end{minipage} & \begin{minipage}[b]{\linewidth}\raggedright
Condizioni
\end{minipage} \\
\midrule\noalign{}
\endhead
\bottomrule\noalign{}
\endlastfoot
\(u^\alpha\) & \(\alpha u^{\alpha-1}u'\) & dove la potenza è reale e derivabile; per \(\alpha\in\mathbb R\) arbitrario, \(u>0\) \\
\(\sqrt{u}\) & \(u'/(2\sqrt{u})\) & \(u>0\) \\
\(\lvert u\rvert\) & \(\operatorname{sgn}(u)u'\) & \(u\ne0\) \\
\(e^u\) & \(e^u u'\) & --- \\
\(a^u\) & \(a^u\ln(a)\,u'\) & \(a>0\) \\
\(\ln\lvert u\rvert\) & \(u'/u\) & \(u\ne0\) \\
\(\sin u\) & \(\cos u\,u'\) & --- \\
\(\cos u\) & \(-\sin u\,u'\) & --- \\
\(\tan u\) & \(u'/\cos^2u\) & \(\cos u\ne0\) \\
\(\arcsin u\) & \(u'/\sqrt{1-u^2}\) & \(\lvert u\rvert<1\) \\
\(\arccos u\) & \(-u'/\sqrt{1-u^2}\) & \(\lvert u\rvert<1\) \\
\(\arctan u\) & \(u'/(1+u^2)\) & --- \\
\(\operatorname{arccot}u\) & \(-u'/(1+u^2)\) & --- \\
\(\cot u\) & \(-u'/\sin^2u\) & \(\sin u\ne0\) \\
\(\sec u\) & \(\sec u\tan u\,u'\) & \(\cos u\ne0\) \\
\(\csc u\) & \(-\csc u\cot u\,u'\) & \(\sin u\ne0\) \\
\(\sinh u\) & \(\cosh u\,u'\) & --- \\
\(\cosh u\) & \(\sinh u\,u'\) & --- \\
\(\tanh u\) & \(u'/\cosh^2u\) & --- \\
\(\coth u\) & \(-u'/\sinh^2u\) & \(u\ne0\) \\
\(\operatorname{arsinh}u\) & \(u'/\sqrt{1+u^2}\) & --- \\
\(\operatorname{arcosh}u\) & \(u'/\sqrt{u^2-1}\) & \(u>1\) \\
\(\operatorname{artanh}u\) & \(u'/(1-u^2)\) & \(\lvert u\rvert<1\) \\
\(\operatorname{arcoth}u\) & \(u'/(1-u^2)\) & \(\lvert u\rvert>1\) \\
\end{longtable}

\section{Derivazione implicita e parametrica}\label{derivazione-implicita-e-parametrica}

Se \(F\in C^1\) in un intorno di \((x_0,y_0)\), \(F(x_0,y_0)=0\) e \(F_y(x_0,y_0)\ne0\), allora localmente \(F(x,y(x))=0\) e

\[
y'(x)=-\frac{F_x(x,y(x))}{F_y(x,y(x))}.
\]

Qui \(F_x\) e \(F_y\) sono le derivate parziali. La condizione \(F_y\ne0\) permette di risolvere localmente l'equazione rispetto a \(y\); se invece \(F_x\ne0\) si può risolvere localmente rispetto a \(x\).

Per una curva parametrica \(x=x(t)\), \(y=y(t)\) con \(x'(t)\ne0\):

\[
\frac{dy}{dx}=\frac{y'(t)}{x'(t)},
\]

\[
\frac{d^2y}{dx^2}
=
\frac{1}{x'(t)}\frac{d}{dt}\left(\frac{y'(t)}{x'(t)}\right).
\]
